\section{\label{sec:summary}Summary}

Quasi and pseudodistributions have opened a new window into nucleon structure and provided an opportunity to extract, for the first time, parton distribution functions directly from lattice QCD. Preliminary results of nonperturbative computations have been encouraging, but there are still challenges to be navigated before the lattice community can claim a robust determination of light-front PDFs. One particular difficulty for lattice QCD is the presence of a power divergence associated with the length of the Wilson line, which must be removed nonperturbatively before one can take a continuum limit. In Ref.~\cite{Monahan:2016bvm}, we introduced the smeared quasidistribution as a tool to circumvent the issue of the power divergence associated with the Wilson line. By taking advantage of the renormalization properties of the gradient flow, a classical evolution of the original degrees of freedom in a new parameter, the flow time, the smeared quasidistribution remains finite in the continuum limit.

Here, I study the smeared quasidistribution at one loop in perturbation theory. Although the gradient flow complicates the perturbative calculation, by introducing new diagrams not present at vanishing flow time and complicating the corresponding integrals, the advantage is significant: provided one fixes the flow time in physical units, the smeared quasidistribution is finite in the continuum limit. The resulting continuum smeared quasidistributions can then be related, directly or through the quasidistribution at zero flow time, to the light-front PDFs in the $\ms$ scheme. The IR behavior of the quasidistribution is unaffected by the flow time, which serves only to regulate the UV behavior, so the matching procedure can be carried out perturbatively, by evaluating the Wilson-line operator between external, gauge-fixed quark states.

I undertake the matching at fixed Wilson-line length, in which case the external quark momentum can be set to zero. This reduces the number of diagrams that contribute at a given order in perturbation theory and simplifies the corresponding one-loop integrals, which can then be carried out analytically. I use dimensional regularization to regularize the logarithmic IR divergences and demonstrate that the resulting poles in the spacetime dimension cancel in the matching relation between the smeared and unsmeared quasidistributions. I show that the limit of vanishing Wilson-line length is well defined and reduces to the known vector current result at finite flow time. In the small flow-time limit, the matrix element depends only logarithmically on the ratio of the Wilson-line length and the gradient flow smearing radius. The smeared matrix element therefore satisfies an equation similar to the typical renormalization-group equation, which provides the potential for defining a nonperturbative step-scaling procedure. Finally, I study the smeared matrix element at finite external momentum by evaluating the corresponding one-loop diagrams numerically. Extending this study to two loops would likely require an automated approach, because the extra vertices induced by the gradient flow cause the number of diagrams to increase rapidly with each order in perturbation theory. A nonperturbative study of the systematic uncertainties associated with the smeared quasi- and pseudodistributions is underway.
